% ABSTRACT

\begin{resumo}[ABSTRACT (OBRIGATÓRIO)]
\begin{Spacing}{1.5}

An essential competency in Computing education is the ability to understand, implement, and select suitable data structures for different problems. In this context, the advance of generative artificial intelligence tools, such as ChatGPT, has raised discussions about their possible effects on students' learning in programming courses. This work investigates the impact of ChatGPT usage on students' learning in the Data Structures and Algorithms course, focusing on Heap and Hash Table topics. The research is methodologically inspired by an experimental study conducted with undergraduate students in data structures, adapting its scope to the topics selected in this work. The proposed study involves practical activities with groups of students under different ChatGPT usage conditions, allowing the comparison of learning evidence through questionnaires, academic performance, and code similarity analysis. Thus, this work seeks to observe whether the use of the tool influences students' conceptual and practical understanding of the topics, as well as whether it contributes to greater similarity among the produced solutions. The expected contribution is to support discussions about the pedagogical use of generative artificial intelligence in Computing courses and the design of assessment activities more suitable for the current educational scenario.

\end{Spacing}

\vspace{1cm}
\textbf{Keywords}: ChatGPT, generative artificial intelligence, data structures and algorithms, Heap, hash tables, learning.

\end{resumo}

% OBSERVAÇÕES:

% Escolha de 3 a 5 palavras ou termos que descrevam bem o seu trabalho .
% As palavras-chave são separadas por vírgula.

